% from ega0-19-fr.tex
\oldpage[0\textsubscript{IV-1}]{108}

où $A_0=A/\mathfrak J=\gr_0(A)$, et $\varphi_M$ et $\varphi_N$ sont les
applications canoniques ($\mathbf 0_{\mathrm{III}}$, 10.1.1.2). Par
hypothèse, $\gr_0(u)$ est bijectif ainsi que $\varphi_M$
($\mathbf 0_{\mathrm{III}}$, 10.2.1), et $\varphi_N$ est surjectif~; on en
déduit d'abord que $\gr_0(u)\otimes1$ est bijectif, puis que $\varphi_N$ est
injectif, donc bijectif, et enfin que $\gr(u)$ est bijectif.
\end{proof}

\begin{lemma}[19.7.1.5]
\label{0.19.7.1.5-fr}
Soient $A$ un anneau noethérien, $\mathfrak J$ un idéal de $A$, $B$, $B'$ deux
$A$-algèbres qui sont des anneaux locaux noethériens, les homomorphismes
$A\to B$, $A\to B'$ étant continus pour la topologie $\mathfrak J$-préadique
sur $A$. On suppose que~:
\begin{enumerate}
  \item[$1^{\mathrm o}$] $B$ et $B'$ sont complets pour les topologies
  $\mathfrak J$-préadiques~;
  \item[$2^{\mathrm o}$] $B$ est une $A$-algèbre formellement lisse~;
  \item[$3^{\mathrm o}$] $B'$ est un $A$-module plat.
\end{enumerate}
Posons $A_0=A/\mathfrak J$, et soit
$u_0:B\otimes_A A_0\to B'\otimes_A A_0$ un $A_0$-isomorphisme~; alors il
existe un $A$-isomorphisme $u:B\to B'$ tel que $u_0=u\otimes1$ (ce qui
entraîne que $B'$ est une $A$-algèbre formellement lisse et $B$ un $A$-module
plat).
\end{lemma}

\begin{proof}
Posons $B_0=B\otimes_A A_0$, $B'_0=B'\otimes_A A_0$. Notons que si
$\mathfrak m$ et $\mathfrak m'$ sont les idéaux maximaux de $B$ et $B'$, les
topologies $\mathfrak J$-préadiques sur $B$ et $B'$ sont séparées puisque
$\mathfrak JB\subset\mathfrak m$ et $\mathfrak JB'\subset\mathfrak m'$~; en
outre, comme $\mathfrak JB'$ est fermé dans $B'$ pour la topologie
$\mathfrak J$-préadique, l'homomorphisme composé
$B\to B_0\xrightarrow{u_0}B'_0$, qui est continu pour les topologies
$\mathfrak J$-préadiques, se factorise en $B\xrightarrow{u}B'\to B'_0$, où
$u$ est un $A$-homomorphisme continu (19.3.10). On a évidemment
$u_0=u\otimes1$, et l'hypothèse que $u_0$ est bijectif entraîne qu'il en est de
même de $u$ en vertu de (19.7.1.4).
\end{proof}

\medskip
\noindent\textbf{(19.7.1.6)}\enspace
\emph{Fin de la démonstration.} --- Pour achever de prouver (19.7.1), il faut
montrer que $a)$ entraîne $b)$~; on sait déjà que $a)$ entraîne que $B_0$ est
une $k$-algèbre formellement lisse (19.3.5, (iii)), donc tout revient à prouver
que $B$ est un $A$-module plat. Il revient au même d'établir que $\widehat B$
est un $\widehat A$-module plat (Bourbaki, \emph{Alg. comm.}, chap.~III, §~5,
n\textsuperscript{o}~4, prop.~4), et l'on sait que $\widehat B$ est une
$\widehat A$-algèbre formellement lisse (19.3.6)~; on peut donc se borner au
cas où $A$ et $B$ sont complets. Comme $B_0$ est une $k$-algèbre formellement
lisse, c'est un anneau régulier (19.6.5) et complet
($\mathbf 0_{\mathrm I}$, 6.3.5)~; appliquant (19.7.1.3), on voit qu'il existe
une $A$-algèbre $B'$ qui est un anneau local noethérien complet et un
$A$-module plat, un homomorphisme local $A\to B'$ et un $k$-isomorphisme
$B_0\simeq B'_0=B'\otimes_A k$. Il suffit alors d'appliquer (19.7.1.5) en
prenant pour $\mathfrak J$ l'idéal maximal de $A$, pour obtenir que $B$ est
$A$-isomorphe à $B'$, donc est un $A$-module plat. C.Q.F.D.

\begin{theorem}[19.7.2]
\label{0.19.7.2-fr}
Soient $A$ un anneau local noethérien, $\mathfrak J$ un idéal contenu dans
l'idéal maximal de $A$, $A_0=A/\mathfrak J$, $B_0$ un anneau local noethérien
complet, $A_0\to B_0$ un homomorphisme local faisant de $B_0$ une
$A_0$-algèbre formellement lisse. Alors il existe un anneau local noethérien
complet $B$, un homomorphisme local $A\to B$ faisant de $B$ un $A$-module
plat, et un $A_0$-isomorphisme
$u:B\otimes_A A_0\xrightarrow{\sim}B_0$. Si $(B',u')$ est un couple
satisfaisant aux mêmes conditions que $(B,u)$, il existe un $A$-isomorphisme
$v:B\xrightarrow{\sim}B'$ rendant commutatif le diagramme
\[
\xymatrix{
  B\otimes_A A_0\ar[rr]^{v\otimes1}_{\sim}\ar[dr]^{\sim}_u &&
  B'\otimes_A A_0\ar[dl]^{u'}_{\sim} \\
  & B_0
}
\]
\end{theorem}

